\documentclass[11pt]{amsart}
\usepackage[utf8]{inputenc}
\usepackage[T1]{fontenc}
\usepackage{amsmath,amssymb,amsthm}
\usepackage[margin=1.1in]{geometry}
\usepackage[hidelinks]{hyperref}
\allowdisplaybreaks
\emergencystretch=2em

\newcommand{\recordurl}{https://github.com/KokunoYumeto/mathematics-workbenches}

\newtheorem{theorem}{Theorem}[section]
\newtheorem{proposition}[theorem]{Proposition}
\newtheorem{lemma}[theorem]{Lemma}
\newtheorem{corollary}[theorem]{Corollary}
\theoremstyle{remark}
\newtheorem{remark}[theorem]{Remark}
\newtheorem{question}[theorem]{Question}

\newcommand{\Fp}{\mathbb{F}_p}
\newcommand{\Zp}{\mathbb{Z}_p}
\newcommand{\Qp}{\mathbb{Q}_p}
\newcommand{\D}{\mathbb{D}}
\newcommand{\Ocris}{\mathcal{O}_{\mathrm{cris}}}
\newcommand{\cris}{\mathrm{cris}}
\newcommand{\Hom}{\operatorname{Hom}}
\newcommand{\Spec}{\operatorname{Spec}}
\newcommand{\mup}{\mu_{p^\infty}}
\newcommand{\QZ}{\Qp/\Zp}

\title[Not full over a reduced base]{Grothendieck's crystalline Dieudonn\'e functor\\ is not full over a reduced base}
\author{}
\date{}

\begin{document}

\begin{abstract}
Grothendieck's crystalline Dieudonn\'e functor sends $p$-divisible groups over a base of characteristic~$p$ to
Dieudonn\'e crystals. Berthelot and Messing showed that it is not full over the non-reduced ring
$\Fp[x,y]/(x^2,xy,y^2)$. De Jong, de Jong--Messing and Lau proved full faithfulness over bases with a $p$-basis, over
excellent local complete intersections and over $F$-finite bases whose divided-power envelopes are $p$-torsion free.
De Jong also proved that over every reduced base of finite type over a field with a finite $p$-basis, every morphism of
Dieudonn\'e crystals comes from a homomorphism up to a torsion term. We show that this torsion term occurs. Let $R$ be
the coordinate ring of two planes in affine $4$-space over $\Fp$ meeting at a point; $R$ is reduced, of finite type and
$F$-finite. Over $R$ there is a nonzero morphism of Dieudonn\'e crystals between $\D(\mup)$ and $\D(\QZ)$, killed
by~$p$, while every homomorphism between $\mup$ and $\QZ$ is zero. So reducedness cannot replace the
complete-intersection or torsion-freeness hypotheses of the positive theorems; the same holds over a reduced
non-noetherian ring. Along the way we correct the formula for divided powers in a standard example at $p=2$. The
examples and proofs were produced by GPT-6 Astra (ChatGPT Pro mode) and checked by Claude Opus~5.5 at effort max; no
human has checked them.
\end{abstract}

\maketitle

\section{Introduction}

Let $p$ be a prime and $S$ a scheme of characteristic~$p$. Grothendieck's crystalline Dieudonn\'e theory, announced in
his lecture at the Nice congress of 1970 and constructed by Berthelot, Breen and Messing~\cite{BBM}, attaches to every
$p$-divisible group $G$ over $S$ a Dieudonn\'e crystal $\D(G)$. This is a finite locally free crystal on the
crystalline site of $S$ over $\Zp$, with Frobenius and Verschiebung $F\colon\varphi^*\D(G)\to\D(G)$ and
$V\colon\D(G)\to\varphi^*\D(G)$ satisfying $FV=VF=p$. The functor $\D$ is contravariant. Whether it is fully faithful
is the basic question of the theory. In his notes for the Nice lecture Grothendieck wrote that the functor ``might be
fully faithful'' (Montpellier archive, cote 129, scan 114, in our reading of the manuscript).

\subsection*{What was known}
De Jong's survey~\cite[pp.~260--261]{dJICM} lists the state of the question. Fully faithful:
\begin{itemize}
\item over perfect fields (classical Dieudonn\'e theory);
\item over schemes that locally have a $p$-basis (Berthelot--Messing~\cite{BM});
\item over excellent schemes whose complete local rings are complete intersections (de Jong--Messing,
  \cite[Thm.~4.6]{dJM}).
\end{itemize}
Fully faithful up to isogeny:
\begin{itemize}
\item over schemes of finite type over a field (de Jong~\cite[Cor.~5.1.2]{dJ95} and~\cite{dJM});
\item over excellent local rings~\cite{dJM}.
\end{itemize}
More recently, Lau~\cite[Cor.~1.7 = Thm.~7.9]{Lau} proved that $\D$ is fully faithful over every $F$-finite scheme
whose divided-power envelopes are $p$-torsion free. He also proved~\cite[Cor.~1.6 = Cor.~9.5]{Lau} that over every
$F$-finite, $F$-nilpotent scheme the map $\Hom(G,H)\to\Hom(\D(H),\D(G))$ is injective with cokernel killed by~$p$.
Anschütz and Le Bras~\cite{ALB} proved full faithfulness over quasisyntomic rings.

The functor is not full in general. Berthelot and Messing~\cite[4.4.1]{BM}, as cited in~\cite[p.~261, (x)]{dJICM} and
\cite[Rem.~3.4]{Lau}, gave the ring $\Fp[x,y]/(x^2,xy,y^2)$. The mechanism is a nonzero $p$-torsion global section of
the crystalline structure sheaf, written out in the Stacks project~\cite[Example~07LJ]{Stacks}.

That ring has nilpotents. For reduced bases de Jong proved more~\cite[Main Thm.~2 = Thm.~5.1.1]{dJ95}: if $S$ is
reduced and of finite type over a field with a finite $p$-basis, then for $p$-divisible groups $G_1,G_2$ over~$S$,
\[
\Hom(\D(G_2),\D(G_1))=(\text{torsion subgroup})\oplus\D(\Hom_S(G_1,G_2)).
\]
He noted that the torsion subgroup is known to vanish when $S$ is a local complete intersection. Zink~\cite[p.~2]{Zink}
pointed out that over reduced rings the Frobenius of a display determines its inverse Verschiebung, so that the display
attached to a crystal carries no extra structure there.

\subsection*{What is proved here}
Put
\[
R=\Fp[a,b,c,d]/(ac,ad,bc,bd)=\Fp[a,b,c,d]/\bigl((a,b)\cap(c,d)\bigr),
\]
the coordinate ring of the planes $\{a=b=0\}$ and $\{c=d=0\}$ in $\mathbb A^4$, which meet only at the origin.

\begin{theorem}\label{thm:main}
For every prime~$p$, the crystalline Dieudonn\'e functor on $p$-divisible groups over $\Spec R$ is not full.
Precisely, there is a nonzero element $\tau$ of $H^0_{\cris}(\Spec R/\Zp)$ with $p\tau=0$ and $\varphi(\tau)=0$.
Multiplication by $\tau$ is a nonzero morphism of Dieudonn\'e crystals $\D(\mup)\to\D(\QZ)$. On the other hand
$\Hom(\QZ,\mup)=0$ and $\Hom(\mup,\QZ)=0$ over~$R$.
\end{theorem}

The ring $R$ is reduced, of finite type over $\Fp$ (hence excellent and $F$-finite) and $F$-nilpotent. It is not
Cohen--Macaulay, hence not a complete intersection: removing the origin disconnects $\Spec R$.

\begin{corollary}\label{cor:main}
\begin{enumerate}
\item The torsion subgroup in de Jong's theorem is nonzero for the reduced scheme $\Spec R$ of finite type over~$\Fp$.
\item The hypotheses ``local complete intersection'' in de Jong--Messing's theorem and ``divided-power envelopes $p$-torsion
  free'' in Lau's theorem cannot be replaced by ``reduced'', not even for $F$-finite schemes of finite type over~$\Fp$.
\item The cokernel in Lau's isogeny theorem, which is killed by~$p$, can be nonzero over a reduced base.
\item The functor is not full over the reduced non-noetherian ring $R[t_1,t_2,\dots]$.
\end{enumerate}
\end{corollary}

This matters as a boundary of the positive theory. All of its integral results for singular bases assume a complete
intersection or a torsion-free envelope, while its isogeny results hold for every reduced base of finite type. One
might have hoped that reducedness bridges the two; the example shows that it does not.

\subsection*{Credit}
The construction of $\tau$ on the non-reduced ring $R_0=\Fp[x,y]/(x^2,xy,y^2)$, its properties $p\tau_0=0$,
$d\tau_0=0$, $\varphi(\tau_0)=0$, and its detection by a finite divided-power algebra are in the Stacks
project~\cite[Example~07LJ]{Stacks}. The failure of fullness over $R_0$ is due to Berthelot and Messing~\cite[4.4.1]{BM}.

Supplied here:
\begin{itemize}
\item the reduced lift $\tau$ on the two planes;
\item the observation that over a reduced base the Hom groups vanish in both directions, which makes the conclusion
  independent of conventions;
\item the non-noetherian variant;
\item a corrected formula for the detector at $p=2$ (Remark~\ref{rem:stacks}). The formula displayed in~\cite{Stacks}
  is correct for odd~$p$, but at $p=2$ it violates the addition axiom of divided powers.
\end{itemize}
The proofs rest on Berthelot's computation of crystalline cohomology by the divided-power de~Rham
complex~\cite[Prop.~07LG]{Stacks}.

\subsection*{Scope} We do not know whether the functor is full over reduced excellent bases that are normal, or
Cohen--Macaulay (Section~\ref{sec:open}).

\section{Proof overview}

To contradict fullness over a reduced base we need a morphism of Dieudonn\'e crystals that comes from no homomorphism.
The simplest groups give the simplest target. Both $\D(\mup)$ and $\D(\QZ)$ are the structure crystal $\Ocris$, with
$(F,V)=(p,1)$ and $(1,p)$ respectively. A morphism between them is multiplication by a global section $\tau$ of
$\Ocris$, and compatibility with $F$ and $V$ reads $p\tau=\varphi(\tau)$. Over a reduced base of characteristic~$p$ the
only $p$-power root of unity is~$1$, so $\Hom(\QZ,\mup)=0$. So it suffices to find a nonzero global section with
$p\tau=\varphi(\tau)$, and a section killed by both $p$ and $\varphi$ will do.

Global sections of $\Ocris$ are the horizontal elements of the $p$-completed divided-power envelope $D$ of a
presentation. Torsion in $D$ comes from relations among products. On two planes the relation is
$(ac)(bd)=(ad)(bc)$. Taking $p$-th divided powers of the two factorizations gives
$\tau=\gamma_p(ac)\gamma_p(bd)-\gamma_p(ad)\gamma_p(bc)$. Multiplying by $p!$ turns both products into
$\gamma_p(abcd)$, so $p\tau=0$. The same identity gives $d\tau=0$. A Frobenius lift multiplies $\tau$ by a multiple
of~$p$, so $\varphi(\tau)=0$.

The difficulty is to see that $\tau\neq0$: $D$ is large and has torsion, so no counting argument applies. The idea is
to map $D$ to a nine-dimensional $\Fp$-algebra $B_p$ with explicitly constructed divided powers, in which $\tau$
becomes a visible basis vector~$uv$. The map sends $a,c\mapsto x$ and $b,d\mapsto y$, so it factors through the
classical ring $R_0$. Reducedness of $R$ is used only for the vanishing of the Hom groups. The specialization is used
only to prove nonvanishing.

\section{Crystalline preliminaries}

Let $S=\Spec A$ with $A$ an $\Fp$-algebra, and work on the crystalline site of $S$ over $\Zp$ with the canonical
divided powers on $(p)$. A morphism of crystals $\Ocris\to\Ocris$ is multiplication by a global section of $\Ocris$.
Write $H^0_\cris(S)$ for these sections and $\varphi$ for the endomorphism induced by absolute Frobenius.

\begin{lemma}\label{lem:hom}
In the contravariant convention, $\D(\mup)=(\Ocris,F=p,V=1)$ and $\D(\QZ)=(\Ocris,F=1,V=p)$. A morphism of Dieudonn\'e
crystals $\D(\mup)\to\D(\QZ)$ is multiplication by some $\tau\in H^0_\cris(S)$ with $p\tau=\varphi(\tau)$, and every such
$\tau$ gives one.
\end{lemma}

\begin{proof}
Both groups are pulled back from $\Spec\Fp$, and $\D$ commutes with base change~\cite{BBM}. Over $\Fp$ the
contravariant Dieudonn\'e module of the \'etale group $\QZ$ has bijective $F$, and that of the multiplicative group
$\mup$ has bijective $V$. In the rank-one case these are the stated scalars. Equivalently, the Hodge filtration is all
of $\D(\mup)_S$ and zero in $\D(\QZ)_S$.

A morphism of crystals $f$ is multiplication by some $\tau$, and $\varphi^*f$ is multiplication by $\varphi(\tau)$.
Compatibility with $F$ says $1\cdot\varphi(\tau)=\tau\cdot p$. Compatibility with $V$ says $\varphi(\tau)\cdot1=p\cdot\tau$.
These are the same equation.
\end{proof}

\begin{remark}\label{rem:convention}
In the other convention, with $F$ and $V$ exchanged on the two groups, the same $\tau$ is a morphism $\D(\QZ)\to\D(\mup)$.
Over a reduced base both $\Hom(\QZ,\mup)$ and $\Hom(\mup,\QZ)$ vanish (Lemma~\ref{lem:groups}), so Theorem~\ref{thm:main}
holds in either convention.
\end{remark}

\begin{lemma}\label{lem:groups}
Let $A$ be a reduced $\Fp$-algebra. Then
\[\Hom_A(\QZ,\mup)=0\quad\text{and}\quad\Hom_A(\mup,\QZ)=0.\]
\end{lemma}

\begin{proof}
A homomorphism $\QZ\to\mup$ is a compatible family of $\zeta_n\in\mu_{p^n}(A)$ with $\zeta_{n+1}^p=\zeta_n$. In
characteristic $p$, $\zeta^{p^n}=1$ means $(\zeta-1)^{p^n}=0$, so $\zeta=1$ in a reduced ring. In the other direction,
$\mu_{p^n}$ is infinitesimal over $\Fp$, and every homomorphism from it to the \'etale group $\mathbb Z/p^n$ is
trivial, over any $\Fp$-algebra.
\end{proof}

Let $P$ be a polynomial ring over $\Zp$, $J\subset P$ an ideal containing $p$, and $A=P/J$. Let $D$ be the $p$-adically
completed divided-power envelope of $(P,J)$, compatible with the divided powers of $p$, and let $d\colon D\to
D\otimes_P\Omega^1_{P/\Zp}$ be its divided-power derivation. Berthelot's theorem~\cite[Prop.~07LG]{Stacks} computes
crystalline cohomology by the divided-power de~Rham complex of $D$. In degree zero:
\begin{equation}\label{eq:H0}
H^0_\cris(\Spec A)=\ker\bigl(d\colon D\to D\otimes_P\Omega^1_{P/\Zp}\bigr).
\end{equation}
If $\sigma\colon P\to P$ is a lift of Frobenius with $\sigma(J)\subseteq J$, then $\sigma$ extends to $D$, and its
restriction to~\eqref{eq:H0} is $\varphi$.

\section{The section \texorpdfstring{$\tau$}{tau}}

Take $P=\Zp[a,b,c,d]$, $J=(p,ac,ad,bc,bd)$, so $P/J=R$. Write $\gamma_n$ for the divided powers in $D$, and put
\[
\tau=\gamma_p(ac)\gamma_p(bd)-\gamma_p(ad)\gamma_p(bc)\in D.
\]

\begin{proposition}\label{prop:tau}
In $D$ we have $p\tau=0$ and $d\tau=0$. For the Frobenius lift $\sigma$ raising each variable to its $p$-th power,
$\sigma(\tau)=0$.
\end{proposition}

\begin{proof}
For $r\in D$ and $s$ in the divided-power ideal, $\gamma_p(rs)=r^p\gamma_p(s)$ and $p!\,\gamma_p(s)=s^p$. Hence
\[
p!\,\gamma_p(ac)\gamma_p(bd)=(ac)^p\gamma_p(bd)=\gamma_p(abcd)=(ad)^p\gamma_p(bc)=p!\,\gamma_p(ad)\gamma_p(bc),
\]
so $p!\,\tau=0$. Since $(p-1)!$ is a unit of $\Zp$, $p\tau=0$.

The derivation satisfies $d\gamma_n(s)=\gamma_{n-1}(s)\,ds$ and $\gamma_{p-1}(s)=s^{p-1}/(p-1)!$. So the coefficient
of $da$ in $d\tau$ is
\[
\frac{a^{p-1}}{(p-1)!}\bigl(c^{p-1}\cdot c\,\gamma_p(bd)-d^{p-1}\cdot d\,\gamma_p(bc)\bigr)
=\frac{a^{p-1}}{(p-1)!}\bigl(\gamma_p(bcd)-\gamma_p(bcd)\bigr)=0.
\]
The coefficients of $db$, $dc$ and $dd$ vanish by the same computation with the letters permuted.

The lift $\sigma$ preserves $J$. For $s$ in the divided-power ideal,
\[
\gamma_p(s^p)=\gamma_p\bigl(p!\,\gamma_p(s)\bigr)=(p!)^p\gamma_p(\gamma_p(s))
=\frac{(p^2)!}{p!}\gamma_{p^2}(s)=(p!)^{p-1}\gamma_p(s)^p.
\]
The last step uses $\gamma_p(s)^p=p!\,\gamma_p(\gamma_p(s))$. Writing $X=\gamma_p(ac)\gamma_p(bd)$ and
$Y=\gamma_p(ad)\gamma_p(bc)$, we get $\sigma(\tau)=(p!)^{2p-2}(X^p-Y^p)$. This is a multiple of $p(X-Y)=p\tau=0$.
\end{proof}

By~\eqref{eq:H0}, $\tau$ is a global section of $\Ocris$ with $p\tau=0=\varphi(\tau)$.

\section{Nonvanishing: a finite divided-power algebra}

Let
\[
B_p=\Fp[x,y,u,v]/(x^3,x^2y,xy^2,y^3,xu,yu,xv,yv,u^2,v^2).
\]
Its monomials $1,x,y,x^2,xy,y^2,u,v,uv$ form an
$\Fp$-basis. Let $K=(x^2,xy,y^2,u,v)$; it has basis $x^2,xy,y^2,u,v,uv$, and $xK=yK=0$, $K^2=\Fp\,uv$, $K^3=0$.

\begin{lemma}\label{lem:detector}
For $z=\alpha x^2+\beta xy+\gamma y^2+\delta u+\varepsilon v+\zeta uv\in K$ define $\gamma_n(z)\in K$ ($n\geq1$) and
$\gamma_0(z)=1$ by
\[
\sum_{n\ge0}\gamma_n(z)T^n=(1+\alpha x^2T+\alpha^puT^p)(1+\beta xyT)(1+\gamma y^2T+\gamma^pvT^p)
(1+\delta uT)(1+\varepsilon vT)(1+\zeta uvT).
\]
This is a divided-power structure on~$K$, with $\gamma_p(x^2)=u$, $\gamma_p(y^2)=v$ and $\gamma_p(xy)=0$. Explicitly,
\[
\sum_n\gamma_n(z)T^n=1+zT+\delta\varepsilon\,uv\,T^2+(\alpha^pu+\gamma^pv)T^p
+(\alpha^p\varepsilon+\gamma^p\delta)\,uv\,T^{p+1}+\alpha^p\gamma^p\,uv\,T^{2p},
\]
where coinciding powers of $T$ are added. In particular $\gamma_2(z)=\alpha^2u+\gamma^2v+\delta\varepsilon\,uv$ when $p=2$.
\end{lemma}

\begin{proof}
The explicit form follows from expanding the product: the only nonzero product of two elements of $K$ is a multiple
of $uv$, coming from $u\cdot v$. In particular $\gamma_n(z)\in K$ for $n\geq1$ and $\gamma_1(z)=z$.

\emph{Addition.} The rule $\gamma_n(z+w)=\sum_i\gamma_i(z)\gamma_{n-i}(w)$ says that the series of $z+w$ is the product
of those of $z$ and~$w$. Each of the six factors is multiplicative in its coefficient. For example
\[
(1+\alpha x^2T+\alpha^puT^p)(1+\alpha'x^2T+\alpha'^puT^p)=1+(\alpha+\alpha')x^2T+(\alpha+\alpha')^puT^p,
\]
because $x^4=x^2u=u^2=0$ and $(\alpha+\alpha')^p=\alpha^p+\alpha'^p$. The other factors are similar, since
$(xy)^2=u^2=v^2=(uv)^2=0$.

\emph{Scalars.} Write $\lambda\in B_p$ as $\lambda_0+\lambda_1x+\lambda_2y+\kappa$ with $\lambda_i\in\Fp$ and $\kappa\in K$.
Then $\lambda z$ has the coordinates $\lambda_0\alpha,\dots,\lambda_0\varepsilon$ and some new $uv$-coordinate. For $n\geq2$,
$\gamma_n(z)$ does not involve $\zeta$ and is homogeneous of degree $n$ in $(\alpha,\dots,\varepsilon)$; hence
$\gamma_n(\lambda z)=\lambda_0^n\gamma_n(z)$. On the other side, $\lambda^n$ acts on $\Fp uv$ as $\lambda_0^n$, and on $u$
and $v$ as $\lambda_0^n$ up to a term $n\lambda_0^{n-1}(\cdot)\,uv$. That term matters only for $n=p$, the only index
where $u$ or $v$ occurs, and then it vanishes. So $\gamma_n(\lambda z)=\lambda^n\gamma_n(z)$.

\emph{Products.} We need $\gamma_i(z)\gamma_j(z)=\binom{i+j}{i}\gamma_{i+j}(z)$. Since $uv\cdot K=0$, the only
products that can be nonzero are
\[
z\cdot z=2\delta\varepsilon\,uv,\qquad z\cdot\gamma_p(z)=(\delta\gamma^p+\varepsilon\alpha^p)uv,\qquad
\gamma_p(z)^2=2\alpha^p\gamma^p\,uv,
\]
and these equal $2\gamma_2(z)$, $(p+1)\gamma_{p+1}(z)$ and $\binom{2p}{p}\gamma_{2p}(z)$, since $\binom{2p}{p}\equiv2\pmod p$.
For $p=2$ the first and third sides are both $0$, and the middle equation is $(\delta\gamma^2+\varepsilon\alpha^2)uv=3\gamma_3(z)$.
For every other pair with $i,j\geq1$ the product vanishes, and $\gamma_{i+j}(z)=0$ or the binomial coefficient is
divisible by~$p$, by Lucas's theorem: $\binom{p}{i}$ and $\binom{2p}{i}$ for $0<i<p$ or $p<i<2p$, and $\binom{p+1}{i}$ for
$1<i<p$.

\emph{Composition.} We need $\gamma_m(\gamma_n(z))=\frac{(mn)!}{m!(n!)^m}\gamma_{mn}(z)$; it is trivial for $m=1$ or
$n=1$. For $m,n\geq2$, $\gamma_n(z)$ lies in $\Fp u+\Fp v+\Fp uv$, with $u$- or $v$-coordinates only for $n=p$. Applying
the definition to $w=\gamma_p(z)$ gives $\gamma_2(w)=\alpha^p\gamma^p\,uv$ and $\gamma_m(w)=0$ for $m\geq3$; every other
$\gamma_m(\gamma_n(z))$ with $m\geq2$ vanishes. The required coefficient of $\gamma_{2p}(z)$ for $(m,n)=(2,p)$ is
$\binom{2p}{p}/2\equiv1\pmod p$, and equals $3\equiv1$ for $p=2$. In the remaining cases with $\gamma_{mn}(z)\neq0$, the
coefficient is divisible by~$p$: these are $mn=p+1$ with $m,n<p$, and $(m,n)=(p,2)$ for odd~$p$.

Finally, $\gamma_p(x^2)=u$ and $\gamma_p(y^2)=v$ read off the first and third factors, and the series of $xy$ is
$1+xyT$.
\end{proof}

\begin{remark}\label{rem:stacks}
The Stacks project~\cite[Example~07LJ]{Stacks} uses the same algebra with the map $\delta_p(z)=\alpha^pu+\gamma^pv$
($K\to K$). For odd $p$ this is the $\gamma_p$ above. For $p=2$ it is not a divided-power operation:
$\gamma_2(u+v)$ must equal $\gamma_2(u)+\gamma_2(v)+uv=uv$, while $\delta_2(u+v)=0$. The correct $\gamma_2$ is given in
Lemma~\ref{lem:detector}. With it, the conclusion of that example holds at $p=2$ as well.
\end{remark}

\begin{proposition}\label{prop:nonzero}
$\tau\neq0$ in $D$.
\end{proposition}

\begin{proof}
The ring map $P\to B_p$ with $a,c\mapsto x$ and $b,d\mapsto y$ kills $p$ and sends $J$ into~$K$, because $ac,ad,bc,bd$
go to $x^2,xy,xy,y^2$. Since $pB_p=0$, the divided powers of Lemma~\ref{lem:detector} are compatible with those of $p$.
By the universal property of the divided-power envelope, the map extends to a divided-power homomorphism $D\to B_p$; it
factors through the $p$-adic completion because $B_p$ is killed by~$p$. It sends $\tau$ to
$\gamma_p(x^2)\gamma_p(y^2)-\gamma_p(xy)^2=uv\neq0$.
\end{proof}

\section{Proof of the theorem and further remarks}

\begin{proof}[Proof of Theorem~\ref{thm:main}]
By Proposition~\ref{prop:tau} and~\eqref{eq:H0}, $\tau$ is a global section of $\Ocris$ on $\Spec R$ with
$p\tau=0=\varphi(\tau)$. It is nonzero by Proposition~\ref{prop:nonzero}. By Lemma~\ref{lem:hom}, multiplication by
$\tau$ is a nonzero morphism of Dieudonn\'e crystals $\D(\mup)\to\D(\QZ)$.

The ideal $(ac,ad,bc,bd)$ equals $(a,b)\cap(c,d)$, an intersection of primes, so $R$ is reduced. By
Lemma~\ref{lem:groups}, $\Hom_R(\QZ,\mup)=0$. Hence the morphism is not $\D$ of any homomorphism, and $\D$ is not
full.
\end{proof}

\begin{proof}[Proof of Corollary~\ref{cor:main}]
(1) The element $\tau$ is killed by $p$, so it lies in the torsion subgroup of the group of morphisms $\D(\mup)\to\D(\QZ)$, while no homomorphism $\QZ\to\mup$ is nonzero.

(2) $R$ is of finite type over $\Fp$, so excellent and $F$-finite. It is not Cohen--Macaulay, since its punctured
spectrum at the origin is disconnected. And $\D$ is not full over it.

(3) $R$ is $F$-finite and $F$-nilpotent, so Lau's theorem applies, and the class of $\tau$ is a nonzero element of the
cokernel.

(4) $R_\infty=R[t_1,t_2,\dots]$ is reduced and not noetherian. Pull $\tau$ back along $R\to R_\infty$. The
retraction $t_i\mapsto0$ shows that the pullback is nonzero, and it still satisfies $p\tau=0=\varphi(\tau)$.
Lemmas~\ref{lem:hom} and~\ref{lem:groups} apply over $R_\infty$.
\end{proof}

\begin{remark}\label{rem:R0}
The same construction on $R_0=\Fp[x,y]/(x^2,xy,y^2)$ gives the section
$\tau_0=\gamma_p(x^2)\gamma_p(y^2)-\gamma_p(xy)^2$ of~\cite[Example~07LJ]{Stacks} and the classical failure of fullness
of Berthelot and Messing~\cite[4.4.1]{BM}. The substitution $a,c\mapsto x$, $b,d\mapsto y$ defines a map $R\to R_0$
carrying $\tau$ to $\tau_0$. In this sense the two planes are a reduced lift of the classical example.
\end{remark}

\begin{remark}\label{rem:witt}
The morphism $\tau$ vanishes on every divided-power thickening $(U,T)$ with $\mathcal O_T$ $p$-torsion free, because its
value there is killed by~$p$. For example, it vanishes on the Witt vector thickenings of a reduced ring, on which the
display theory of Zink and Lau lives. So $\tau$ is invisible to displays. This is consistent with Zink's
remark~\cite[p.~2]{Zink} that over reduced rings the inverse Verschiebung is recovered from the Frobenius. The failure
of fullness lives in the torsion of the crystalline site, as de Jong's decomposition predicts.
\end{remark}

\section{What remains open}\label{sec:open}

The example is reduced, and even seminormal, but it is not $S_2$: its two components meet in codimension two. We do
not know the answer to the following.

\begin{question}
Is the crystalline Dieudonn\'e functor full over every reduced excellent $\Fp$-scheme that is normal? Over every one
that is Cohen--Macaulay?
\end{question}

By Lau's theorem a positive answer would follow from $p$-torsion freeness of the divided-power envelopes; we do not know
whether normal rings of finite type over $\Fp$ have this property. The example also leaves open a classification of the
torsion subgroup in de Jong's theorem. Over $R$ we have exhibited one nonzero element, killed by $p$; we have not
computed the full cokernel.

\section*{A note on method}
The mathematics of this note was produced by GPT-6 Astra (OpenAI, ChatGPT Pro mode), in a research run on
Grothendieck's manuscripts on crystalline Dieudonn\'e theory. The run found the reduced lift, the non-noetherian
variant and the corrected divided-power formula at $p=2$.

Claude Opus~5.5 at effort max (Anthropic) checked the claims:
\begin{itemize}
\item targeted proof checks of $p\tau=0$, $d\tau=0$, $\varphi(\tau)=0$ and of the vanishing Hom groups;
\item an independent verification by hand of every divided-power axiom in Lemma~\ref{lem:detector};
\item the literature comparison above, reading at the source de Jong~\cite{dJ95,dJICM}, de Jong--Messing~\cite{dJM},
  Lau~\cite{Lau}, Zink~\cite{Zink} and the Stacks example~\cite{Stacks}.
\end{itemize}
The statements of~\cite{BBM,BM,ALB} are cited through those sources and from memory; Berthelot--Messing~\cite{BM} could
not be consulted. Claude Opus~5.5 at effort max also edited the text. No human has checked this note.

Our text is dedicated to the public domain under CC0~1.0. The records of the checks are at \url{\recordurl}.

\begin{thebibliography}{99}

\bibitem{ALB} J.~Ansch\"utz and A.-C.~Le~Bras, \emph{Prismatic Dieudonn\'e theory}, Forum Math. Pi \textbf{11} (2023),
e2; arXiv:1907.10525.

\bibitem{BBM} P.~Berthelot, L.~Breen and W.~Messing, \emph{Th\'eorie de Dieudonn\'e cristalline II}, Lecture Notes in
Math.~930, Springer, 1982.

\bibitem{BM} P.~Berthelot and W.~Messing, \emph{Th\'eorie de Dieudonn\'e cristalline III: th\'eor\`emes d'\'equivalence et
de pleine fid\'elit\'e}, in: The Grothendieck Festschrift I, Progr. Math.~86, Birkh\"auser, 1990, 173--247.

\bibitem{dJ95} A.~J.~de~Jong, \emph{Crystalline Dieudonn\'e module theory via formal and rigid geometry},
Publ. Math. IH\'ES \textbf{82} (1995), 5--96.

\bibitem{dJICM} A.~J.~de~Jong, \emph{Barsotti--Tate groups and crystals}, Doc. Math., Extra Vol. ICM 1998, II, 259--265.

\bibitem{dJM} A.~J.~de~Jong and W.~Messing, \emph{Crystalline Dieudonn\'e theory over excellent schemes},
Bull. Soc. Math. France \textbf{127} (1999), 333--348.

\bibitem{Lau} E.~Lau, \emph{Divided Dieudonn\'e crystals}, arXiv:1811.09439v2.

\bibitem{Stacks} The Stacks Project Authors, \emph{The Stacks project}, Section 60.22, Example~60.22.1 (Tag 07LJ), and
Proposition~60.21.3 (Tag 07LG), \url{https://stacks.math.columbia.edu}.

\bibitem{Zink} T.~Zink, \emph{A Dieudonn\'e theory for $p$-divisible groups}, in: Class field theory---its centenary
and prospect, Adv. Stud. Pure Math.~30, 2001, 139--160.

\end{thebibliography}

\end{document}
